\documentclass[11pt,a4paper]{article}
\usepackage[utf8]{inputenc}
\usepackage[T1]{fontenc}
\usepackage[french]{babel}
\usepackage[margin=2.5cm]{geometry}
\usepackage{enumitem}
\usepackage{hyperref}

\title{Compte Rendu 4}
\author{Matis Bégot, Rémy Guibert, Lucas Paulo}
\date{15 novembre 2026}

\begin{document}
\maketitle

\subsection{Avancement modèle}
On a commencé par entraîner un modèle sur une seule caractéristique ('Pouces distingués'), il y a 1/3 d'image avec le label 'oui' et 2/3 'non' (sur 150 images), ce qui est relativement équilibré. On obtient ces résultats avec des batches de taille 16 et 8 epochs :

![image](image.png)

Note : Il y a de "l'overfitting": le modèle fonction bien mieux sur les données d'entrainement par rapport au données de test.

\subsection{Potentiel approfondissement}
Il existe aussi la méthode Zero-shot Learning. Cette méthode aurait un intérêt pour reconnaître des classes non vues, des classes / labels sur lesquels notre modèle n'aurait pas été entraîné.

\section{Configuration de l'environnement}
Nous avont utilisé l'outil \verb|uv| afin de gérer facilement Python, sa version et ses dépendances avec une installation facile\section{Autre techniques d'apprentissage}
Nous avons commencé à nous renseigner et à comprendre comment fonctionne l'apprentissage Few-Shot Learning.
Cette méthode nous permettrait d'avoir un entraînement plus adapté à la taille de notre base de données.

\end{document}